\section{An elimination belongs to an identified transaction}
\label{sec:group-elimination-joins}

Equal balances need not describe the same debt. Suppose the parent records
a receivable of 100 from subsidiary S, while another subsidiary T reports a
payable of 100. Cancelling those numbers would balance the worksheet while
misstating both relationships. The group must first establish the seller,
buyer and economic event that connect the two entries. IFRS~10.B86 requires
the elimination of intragroup transactions; it does not make equality of
amounts evidence that two entries are counterparts
\citep[IFRS~10.B86--88, B92--93]{eu20231803}.

For each proposed elimination, retain the reporting group and perimeter
version, both legal entities and book revisions, the common economic-event
reference, invoice or settlement allocations, transaction currency, reporting
cutoff and history branch. The two books' local event identifiers can differ;
an evidenced mapping connects them to the common transaction. For balance elimination,
use the perimeter at the reporting date; for income and cash-flow elimination,
use the intervals during which control exists. A receivable created before
control began can be intragroup at the reporting date without making the
original sale an intragroup revenue transaction for that period. Acquisition
measurement and subsequent movements must remain distinct. Align accounting
policies, reporting dates and currency measurements before comparing balances. A reporting-date exception
needs the conditions and intervening adjustments in IFRS~10.B92--93.

One invoice may correspond to several ledger lines. An explicit allocation
of 40 and 60 can exhaust a receivable of 100; reusing the 40 creates an
invalid allocation of 140. Each source amount is consumed at most once in a
given consolidation version. A receivable of 100 against a payable of 90
leaves an unresolved difference of 10 until a settlement, cutoff correction
or other supported adjustment explains it. Equal labels, a balancing entry
or a stale source revision cannot supply that explanation. Unknown matches
remain visible and prevent a complete-consolidation result for that view.
These identity rules are implementation choices for preserving the source
meaning; they are not a claim that IFRS prescribes a database key.

Now follow one correctly matched inventory sale. The wholly owned subsidiary
begins with cash and capital of 50. The parent holds cash 90, inventory 60 and
the investment 50, with equity 200. All amounts use one currency; taxes,
impairment and other transactions are absent by assumption. The parent sells
the inventory to the subsidiary for 100 on credit. The subsidiary then sells
40 percent of those goods to an external customer for cash 60, recording
cost 40 and retaining inventory 60. The parent's profit is 40 and the
subsidiary's profit 20, but the group has earned only 36: external revenue 60
less the original cost 24 of the goods sold.

Eliminate the investment 50 against subsidiary opening equity 50 and the
payable 100 against receivable 100. Of the internal markup 40, the fraction
still in inventory is 60 percent. The performance adjustment therefore debits
internal revenue 100 and credits cost of sales 76 and inventory 24. Combined
cost of sales 100 becomes 24, inventory 60 becomes 36, and profit 60 becomes 36.
Closing group cash 200 plus inventory 36 equals opening equity 200 plus
profit 36. The same result follows directly from the external history:
opening group cash 140 plus customer receipt 60, with inventory cost 60 less 24.
The adjustment retains the lot and external-sale references needed to release
the remaining markup when those goods leave the group. Actual tax bases,
upstream NCI allocation and evidence that an intragroup loss reflects impairment
require their own consequences under IFRS~10.B86(c). The earlier foreign-loan
case also remains: eliminating a monetary balance does not automatically erase
its exchange gain \citep[IFRS~10.B86(c); IAS~21.45]{eu20231803}.

\section{Restating purchasing power before translating currency}
\label{sec:group-hyperinflation}

An inventory cost and a bank balance of the same nominal amount do not retain
the same purchasing power. IAS~29 addresses this problem when the functional
currency belongs to a hyperinflationary economy. Its paragraph 3 requires
judgment across economic indicators; three-year inflation approaching or
exceeding 100 percent is one indicator, not an automatic test. The selected
assessment, effective period and general price index must therefore accompany
the entity book \citep[IAS~29.1--4, 8--19, 24--28, 37--39]{eu20231803}.

For historical-cost nonmonetary amounts, retain each acquisition or measurement
date and restate from that date's index to the closing index. Restate income
and expense from recognition dates, and contributed equity from the relevant
contribution dates. Monetary balances remain at their closing nominal amounts;
contractually index-linked amounts first follow their contractual adjustment.
Amounts already measured at closing current values are not indexed again.
After restatement, apply the relevant recoverability or net-realizable-value
limit. Missing dates or indices require the specific permitted estimation
procedure and its support; an unexplained common multiplier is insufficient.

Consider a period beginning with cash 100, inventory 60, fixed nominal debt 80
and contributed capital 80. All arose at the opening date, and there are no
transactions, interest, tax or impairment during the period. Assume IAS~29
already applies and the general index rises from 100 to 150. Inventory becomes 90
and capital 120; cash 100 and debt 80 remain unchanged. Holding cash loses 50
units of closing purchasing power, while owing the fixed debt gains 40.
The net monetary loss is therefore 10. If \(M\) denotes the constant excess
of monetary assets over monetary liabilities, and \(I_0,I_1\) the opening
and closing indices, this special case gives
\[
 L_{\mathrm{mon}}=M\left(\frac{I_1}{I_0}-1\right)
 =(100-80)(1.5-1)=10.
\]
Positive \(L_{\mathrm{mon}}\) is a loss. A net monetary liability reverses
the sign; an unchanged price index gives zero. The restated balance sheet
reconciles as \(100+90-80=110=120-10\). The restatement entry debits
inventory 30 and monetary loss 10 and credits restated capital 40.
If inventory's supported closing net realizable value is 84, a separate
write-down 6 reduces equity to 104. Indexing both monetary balances would
incorrectly increase net assets by 10 and conceal the purchasing-power loss.

The constant-exposure formula does not apply to arbitrary turnover.
For a piecewise constant nominal exposure, split the history at every monetary
change. If \(M_j\) is the net monetary amount held between index observations
\(I_{j-1}\) and \(I_j\), its loss expressed at closing index \(I_T\) is
\(M_j(I_j/I_{j-1}-1)I_T/I_j\); sum those interval losses and separately
handle contractually indexed balances. Recognition times and index observations
determine those intervals. IAS~29.27 also permits a weighted-average estimate;
its approximation must be identified. Reconcile the result with restated
nonmonetary amounts, equity and income rather than inventing a residual profit.
For example, external cash income 30 at index 125 changes the preceding net
monetary exposure from 20 to 50. The two losses are 6 and 10, totaling 16;
restated income 36 and monetary loss 16 reconcile closing net assets 140 to
restated opening capital 120.
Cash-flow items likewise use closing purchasing-power units. In the no-flow
case, opening cash expressed in closing units is 150 and ending cash 100;
the purchasing-power reconciliation of minus 50 is not a payment
\citep[IAS~29.26--34]{eu20231803}.

First application requires earlier acquisition history. IFRIC~7.3--5
requires restatement as though the economy had always been hyperinflationary
and a particular ordering for opening deferred tax. In a separate
inventory example, cost 40 at index 80 becomes 50 at the opening index 100
and 75 at the closing index 150. Assume a fixed nominal tax base 40, tax
rate 20 percent, profit-origin differences and no recognition exception.
Opening deferred tax is \((50-40)\times0.2=2\), then restated to 3.
Closing deferred tax is \((75-40)\times0.2=7\); the period expense is 4.
Indexing a nominal opening deferred-tax balance of zero would miss this
sequence \citep[IFRIC~7.3--5; IAS~29.32]{eu20231803}.

Only then translate the restated statements. At a supplied closing rate of
USD 0.20 per local unit, the first example gives cash 20, inventory 18, debt 16,
capital 24 and monetary loss 2: net assets 22 equal 24 minus 2. Translating
unrestated inventory 60 would report 12 and understate the asset by 6.
IAS~21.42--43 uses the closing rate for these restated current-period amounts.
When the presentation currency is not hyperinflationary, prior comparative
amounts remain those presented in the relevant prior-year statements. Thus
prior cash 100 translated at the prior rate 0.25 remains USD 25, even though
the current local comparative is restated. Keep these two comparative
representations distinct \citep[IAS~21.42--47; IAS~29.34--36]{eu20231803}.

\section{A quoted rate may not be available for the required purpose}
\label{sec:group-exchangeability}

A subsidiary may obtain foreign currency to import goods while being unable
to remit its net investment. Exchangeability therefore belongs to a currency
direction, measurement date and purpose. The inspected EU amendment 2024/2862
requires application no later than financial years beginning on or after
1 January 2025; IAS~21.60L permits earlier application with disclosure.
The following method uses that amendment, rather than extending the older
paragraph 26 rule about temporary lack of exchangeability to later periods
\citep[Article~2; IAS~21.8--8B, 19A, 26, 60L--M]{eu20242862}.

Record the entity, both currencies, measurement date, purpose, amount needed,
amount obtainable, permitted mechanisms and their enforceability, normal
administrative delay and evidence for the finding. The assessment concerns
ability, including indirect exchange, rather than management's intention.
For individual transactions, the purpose is realization or settlement of
those transactions; for presentation-currency translation it is realization
or settlement of net assets or liabilities; for a foreign operation it is
realization or settlement of the net investment. Assess the relevant direction
and purpose separately. Obtaining only an insignificant amount fails the
exchangeability test. The ratio 5 available to 1000 needed is evidence for a
judgment, not a threshold supplied by this design
\citep[IAS~21.A2--A10]{eu20242862}.

When exchangeability fails, IAS~21.19A requires an estimate reflecting an
orderly exchange between market participants at the measurement date under
prevailing economic conditions. A rate observable for another purpose, or
the first subsequent exchange rate, can be used without adjustment only
if it meets that objective. Examine incentives or penalties, restricted
purposes, interventions, update frequency, elapsed time and inflation.
Another estimation technique may adjust observable rates, including those
from mechanisms that would not establish exchangeability. It must preserve
the inputs, adjustments and reasoning supporting the measurement-date
objective. Absence of that support leaves the rate unresolved; the design
does not manufacture a universal estimator
\citep[IAS~21.A11--A17]{eu20242862}.

For a conditional arithmetic illustration, assume an admitted estimate of
USD 0.08 per local unit for a foreign operation's net investment. Local net
assets 1000 translate to 80. A preferential import rate 0.10 would produce 100;
a later rate 0.09 following changed conditions would produce 90. Neither
alternative is interchangeable with the admitted estimate merely because it
is observable. This calculation exposes the consequences of the choice;
the numerical differences do not prove that 0.08 meets paragraph 19A.

Initial application has its own journal origin. Suppose an affected monetary
payable of 100 foreign units moves from 10 to 12 functional units per foreign
unit at initial application. The liability increase 200 adjusts opening
retained earnings under paragraph 60M(a). In a separate nonhyperinflationary
foreign-operation translation example, changing the rate for net assets 1000
from 0.10 to 0.08 reduces translated equity 20; paragraph 60M(b) assigns the
initial effect to accumulated translation differences. Hyperinflationary
functional-currency equity items also use the estimated initial rate.
Comparatives are not restated for this amendment. Keep these transition
adjustments distinct from subsequent-period profit or OCI
\citep[IAS~21.60M]{eu20242862}.

The accompanying disclosures identify affected currencies and restrictions,
transactions and carrying amounts, rates, estimation methods and assumptions,
and exposure to risk. A foreign operation also needs the additional
information in paragraph A20, including relevant support arrangements
\citep[IAS~21.57A--B, A18--A20]{eu20242862}.
These methods extend the documentary group design. They do not silently add
fields or supported events to the closed \texttt{GroupForeignOperation/1}
exchange profile; its successor must explicitly bind the additional histories
and assessments before an implementation can claim these operations.
